% Part: proof-theory
% Chapter: proof-search
% Section: completeness

\documentclass[../../../include/open-logic-section]{subfiles}

\begin{document}

\olfileid{pt}{ps}{cpl}

\olsection{Completeness}

We now show that the proof search algorithm is fool-proof, i.e., if it
aborts or never halts, there is no !!{proof} of the starting
sequent~$\Gamma \Sequent \Delta$. To do this, we define
!!a{structure}~$\Struct M$ such that for all $!A \in \Gamma$,
$\Sat{M}{!A}$, and for all~$! \in \Delta$, $\Sat/{M}{!A}$. In other
words, all !!{formula}s in~$\Gamma$ are true and all !!{formula}s
in~$\Delta$ are false in~$\Struct M$, i.e., $\Gamma \Sequent \Delta$
is not valid. Since \Log{G3c} is sound, $\Gamma \Sequent \Delta$ has
no !!{proof}

We define $\Struct{M}$ from a pair $\Theta, \Xi$ of !!{formula}s and a
set~$C$ of !!{constant}s. We obtain $\Theta, \Xi$ and~$C$ from a
\emph{failure branch} of the failed (aborted or non-terminating) run
of the proof search. A failure branch is a sequence $\Pi_n \Sequent
\Lambda_n$ of sequents and sets of !!{constant}s~$C_n$ such that
$\Pi_0 \Sequent \Lambda_0$ is the end-sequent $\Gamma \Sequent
\Delta$, and $\Pi_{n+1} \Sequent \Lambda_{n+1}$ is a premise of $\Pi_n
\Sequent \Lambda_n$ generated at stage~$n+1$ for which the algorithm
does not find !!a{proof}. Such a premise must obviously exist for
each~$n$, since otherwise the algorithm would have found !!a{proof}.
We let $C_n$ be the set of !!{constant}s associated with~$\Pi_n
\Sequent \Delta_n$.

If the algorithm aborts at a top-most sequent consisting only of
atomic !!{formula}s, the branch ending in that sequent is a failure
branch. If if the algorithm never terminates, it keeps generating
larger and larger trees. You can think of these as growing an infinite
tree (the tree you'd have after infinitely many steps of the algorithm
had run). This would be an infinite tree which is, however, finitely
branching (every node only has one or two children---the sequents
immediately above it). Every infinite, finitely branching tree must
have an infinite branch by K\H{o}nig's Lemma. Such an infinite branch
would be a failure branch in the case where the search algorithm runs
forever.

If the sequence of $\Pi_n \Sequent \Lambda_n$ and $C_n$ is a failure
branch, let $\Theta = \bigcup_n \Pi_n$ and $\Xi = \bigcup_n \Lambda_n$
(removing indices and multiple occurrences of !!{formula}s) and $C =
\bigcup_n C_n$.

We are ready to define~$\Struct{M}$.
\begin{enumerate}
\item The domain $\Domain{M}$ is the
set of terms that can be formed from~$C$ and the !!{function}s
in~$\Gamma \Sequent \Delta$, if there are any, but no !!{variable}s. 
\item If $c \in \Domain{M}$, $\Assign{c}{M} = c$.
\item If $f$ is an $m$-place !!{function} in $\Gamma \Sequent \Delta$, and $t_1$,
\dots, $t_m \in \Domain{M}$, then $\Assign{f}{M}(t_1, \dots, t_m) =
f(t_1, \dots, t_n)$.
\item If $R$ is an $m$-place !!{predicate} in $\Gamma \Sequent
\Delta$, then $\Assign{R}{M} = \Setabs{\tuple{t_1, \dots, t_m} \in
\Domain{M}^n}{R(t_1, \dots, t_n) \in \Theta}$.
\end{enumerate}
$\Struct{M}$ is a \emph{term model} in that its domain is a set of
terms, and the !!{constant}s and !!{function}s are interpreted so as
to guarantee that $\Value{t}{M} = t$. The atomic !!{formula}s
in~$\Theta$ are used to define $\Assign{R}{M}$ in such a way as to
guarantee that $\Sat{M}{R(t_1, \dots, t_n)}$ iff $R(t_1, \dots, t_n)
\in \Theta$. 

\begin{lem}\ollabel{lem:truth}
For all $!A \in \Theta \cup \Xi$:
\begin{enumerate}
  \item If $!A \in \Theta$ then $\Sat{M}{!A}$.
  \item If $!A \in \Xi$ then $\Sat/{M}{!A}$.
\end{enumerate}
\end{lem}

\begin{proof}
  By induction on $\depth{!A}$.

First assume $!A$ is atomic, i.e., either $!A \ident \lfalse$ or $!A
\ident R(t_1, \dots, t_m)$.

Because $\Pi_n \Sequent \Lambda_n$ is a failure branch, no $\Pi_n$ can
contain $\lfalse$ (otherwise $\Pi_n \Sequent \Lambda_n$ would be an
axiom). $\Sat{M}{R(t_1, \dots, t_m)}$ iff $R(t_1, \dots, t_m) \in
\Theta$ by definition of~$\Struct{M}$. Together we have: if $!A \in
\Theta$, then $\Sat{M}{!A}$.

If $!A$ is atomic and $!A \in \Pi_n$ then $!A \in \Pi_{n+1}$ and if
$!A \in \Lambda_n$ then $!A \in \Lambda_{n+1}$. (This follows from the
way the proof search algorithm generates successor sequents.) Hence,
if $!A \in \Theta \cap \Xi$ then for some~$n$, $!A \in \Pi_n \cap
\Lambda_n$. This means that $\Pi_n \Sequent \Lambda_n$ is an axiom,
and the failure branch can contain no axioms. So $!A \notin \Theta
\cap \Xi$ for atomic~$!A$ by construction, and consequently if $!A \in
\Xi$ then $!A \notin \Theta$. This implies that if $!A \in \Xi$ then
$\Sat/{M}{!A}$ by definition of~$\Struct{M}$.

For the inductive step, assume (1) and (2) hold for all !!{formula}s
of lower !!{depth} than~$!A$. We prove (1) and (2) for $!A$ by
distinguishing cases.

First note that when $!A^i$ occurs in $\Pi_n \Sequent \Lambda_n$ then
$!A^i$ also occurs in $\Pi_{n+1} \Sequent \Lambda_{n+1}$ unless $i$ is
the smallest index in $\Pi_n \Sequent \Lambda_n$. Since the smallest
index increases in topmost sequents added at each stage, eventually we
reach $\Pi_n \Sequent \Lambda_n$ in which $!A^i$ occurs, and $i$ is
the smallest index. At stage $n+1$ the algorithm reduces~$!A^i$. In
other words, if $!A \in \Theta$, then for some $n$, $\Pi_n = \Pi_n',
!A^i$ and $i$ is the smallest index. And if $!A \in \Xi$, then for
some $n$, $\Lambda_n = \Lambda_n', !A^i$ and $i$~is the smallest
index.

\begin{enumerate}
  \item $!A \in \Theta$ and $!A \ident !B \land !C$. Then for some
  $n$, $\Pi_n = \Pi_n', (!B \land !C)^i$. $\Pi_{n+1} \Sequent
  \Lambda_{n+1}$ is the corresponding premise of~\LeftR{\land}, i.e.,
  $\Pi_{n+1} = \Pi_n', !B^k, !C^{k+1}$. Consequently, $!B, !C \in
  \Theta$. By inductive hypothesis, $\Sat{M}{!B}$ and $\Sat{M}{!C}$.
  By definition of $\Sat{M}{}$, we have $\Sat{M}{!B \land !C}$.

  \item $!A \in \Xi$ and $!A \ident !B \land !C$. Then for some $n$,
  $\Lambda_n = \Lambda_n', !B \land !C$. $\Pi_{n+1} \Sequent
  \Lambda_{n+1}$ is one of the two premises of~\RightR{\land} with
  conclusion $\Pi_n \Sequent \Lambda'_n, !B \land !C$, i.e.,
  $\Lambda_{n+1} = \Lambda_n', !B^k$ or $\Lambda_{n+1} = \Lambda_n',
  !C^k$. Consequently, either $!B \in \Xi$ or $!C \in \Xi$. By
  inductive hypothesis, either $\Sat/{M}{!B}$ or $\Sat/{M}{!C}$. By
  definition of $\Sat{M}{}$, we have $\Sat/{M}{!B \land !C}$.
  
  \item The cases where $!A \ident \lnot !B$, $!A \ident !B \lor !C$,
  and $!A \ident !B \lif !C$ are similar and left as exercises.

  \item $!A \in \Theta$ and $!A \ident \lexists[x][!B(x)]$. Then for
  some $n$, $\Pi_n = \Pi_n', \lexists[x][!B(x)]^i$. $\Pi_{n+1}
  \Sequent \Lambda_{n+1}$ is the corresponding premise
  of~\LeftR{\lexists}, i.e., $\Pi_{n+1} = \Pi_n', !B(c)^k$ for
  some~$c$. Consequently, $!B(c) \in \Theta$ and~$c \in C$. By
  inductive hypothesis, $\Sat{M}{!B(c)}$. By definition of
  $\Sat{M}{}$, we have $\Sat{M}{\lexists[x][!B(x)]}$.

  \item $!A \in \Xi$ and $!A \ident \lexists[x][!B(x)]$. Then for some
  $n$, $\Lambda_n = \Lambda_n', \lexists[x][!B(x)]^i$. Since every time
  $\lexists[x][!B(x)]$ is reduced, $\lexists[x][!B(x)]$ remains on the
  right side of the premise sequent with a new index,
  $\lexists[x][!B(x)]$ is reduced infinitely often on the right side
  in a failure branch if it occurs there. Every term $t \in
  \Domain{M}$ eventually is ``the first term only containing constants
  in~$C_n$ not already used in a reduction of $\lexists[x][!B(x)]$ on
  the right'' on a failure branch. Hence, for all $t \in \Domain{M}$,
  $!B(t) \in \Lambda_n$ for some~$n$ and consequently $!B(t) \in \Xi$.
  By inductive hypothesis, $\Sat/{M}{!B(t)}$ for all $t \in
  \Domain{M}$. By definition of $\Sat{M}{}$, we have
  $\Sat/{M}{\lexists[x][!B(x)]}$.

  \item The cases where $!A \ident \lforall[x][!B(x)]$ are similar and
  left as exercises.
\end{enumerate}
\end{proof}

\begin{cor}\ollabel{cor:complete} The proof search algorithm for
\Log{G3c} is complete: if it aborts or never halts, the starting
sequent~$\Gamma \Sequent \Delta$ is not valid and hence has no
!!{proof}.
\end{cor}

\begin{proof}
  If the algorithm aborts or never halts, there is a failure branch
  from which $\Struct{M}$ can be defined. By \olref{lem:truth}, for
  all $!A \in \Gamma$, $\Sat{M}{!A}$ and for all $!A \in \Delta$,
  $\Sat/{M}{!A}$. (Recall that $\Pi_0 = \Gamma$ and $\Lambda_0 =
  \Delta$, so $\Gamma \subseteq \Theta$ and $\Delta \subseteq \Xi$.)
  So $\Gamma \Sequent \Delta$ is not valid. By soundness of~\Log{G3c},
  $\Gamma \Sequent \Delta$ has no !!{proof}.
\end{proof}

\begin{cor}
\Log{G3c} is complete: if $\Gamma \Sequent \Delta$ is valid,
$\Log{G3c} \Proves \Gamma \Sequent \Delta$.
\end{cor}

\begin{proof}
  We prove the contrapositive. Suppose $\Log{G3c} \Proves/ \Gamma
  \Sequent \Delta$. Then the proof search algorithm must abort or
  never halt, as there is no !!{proof} to be found. By \olref{cor:complete}, $\Gamma \Sequent \Delta$ is not
  valid.
\end{proof}


\end{document}
